\section{Metody a nástroje pro vizualizaci dat}

\subsection*{Smysl prezentace a vizualizace informací}
Dodaný materiál k vizualizaci dat není systematický, ale obsahuje části o prezentaci informací, tabulkách, grafech, obrázcích, Paretově grafu a histogramu. Pro přípravu lze využít základní myšlenku: vizualizace slouží k tomu, aby data byla čitelná, srozumitelná a využitelná pro rozhodování.

Materiály zdůrazňují vztah obsahu a formy. U prezentace informací nestačí mít správná data; je nutné je podat tak, aby publikum pochopilo hlavní sdělení. Úspěšná prezentace vyžaduje znát publikum, vyjadřovat se stručně, držet se tématu a věřit tomu, co říkám. U datové vizualizace to znamená přizpůsobit graf cílové skupině: jinak bude vypadat manažerský dashboard, jinak analytický graf pro datového specialistu a jinak graf v závěrečné práci.

\subsection*{Tabulky, grafy a obrázky}
Materiály popisují \term{tabulky} jako vhodné pro shrnutí číselných dat nebo přehled různých charakteristik. Každá tabulka má mít číslo, název a zdroj. V datové analytice je tabulka vhodná pro přesné hodnoty, porovnání kategorií a kontrolu detailů. Není ale vždy nejlepší pro rychlé odhalení trendu nebo extrému.

\term{Grafy} slouží jako ilustrace a vysvětlení v textu. Graf má mít popsané osy, legendu, název a zdroj. Graf by měl odpovídat typu dat a otázce, kterou řešíme. U časového vývoje dává smysl liniový graf, u kategorií sloupcový graf, u rozdělení histogram, u vztahu dvou metrických proměnných bodový graf.

\term{Obrázky a náčrty} mají být jednoduché a srozumitelné. Do textu se mají vkládat jen tehdy, pokud podporují výklad. Stejná zásada platí pro dashboardy: vizuální prvek má sloužit sdělení, ne být dekorací.

\subsection*{Paretův graf a histogram}
Materiály uvádějí dva konkrétní typy grafů.

\term{Paretův graf} kombinuje sloupcový a čárový graf. Sloupce ukazují četnosti kategorií seřazené od největší po nejmenší a čára zobrazuje kumulativní četnost v procentech. Hodí se pro hledání hlavních příčin problému nebo pro pravidlo 80/20: malý počet kategorií často vysvětluje velkou část jevu.

\term{Histogram} slouží k zobrazení rozdělení hodnot, typicky rozdělením číselné proměnné do intervalů. Umožňuje vidět tvar rozdělení, šikmost, koncentraci hodnot a odlehlé oblasti.

\subsection*{Vizualizace v BI}
V části o BI materiály vysvětlují, že data mají být manažerům reprezentována v čitelné podobě a že BI poskytuje multidimenzionální pohled na data. Vizualizace je tedy výstupní vrstva BI: reporty, dashboardy, OLAP analýzy a ad hoc přehledy. Důležité je napojení na dimenze a metriky. Například tržby lze zobrazit podle času, regionu, produktu nebo zákaznického segmentu.

\todohead
\begin{itemize}
  \item Doplnit typy grafů podle účelu: porovnání, trend, rozdělení, vztah, kompozice, mapa, hierarchie.
  \item Doplnit zásady dobré vizualizace: volba měřítka, práce s barvou, čitelnost, odstranění zbytečného šumu, správné osy.
  \item Doplnit nástroje: Excel, Power BI, Tableau, Qlik, Python matplotlib, seaborn a plotly, R ggplot2.
  \item Doplnit dashboardy, interaktivitu, drill-down a storytelling s daty.
\end{itemize}

\newpage
